\documentclass[12pt]{article}

\input{../../_assets/preambulo.tex}


\begin{document}

    % 1. Foto de fondo
    % 2. Título
    % 3. Encabezado Izquierdo
    % 4. Color de fondo
    % 5. Coord x del titulo
    % 6. Coord y del titulo
    % 7. Fecha

    
    \input{../../_assets/portada}
    \portadaExamen{ffccA4.jpg}{Álgebra II\\Examen VIII}{Álgebra II. Examen VIII}{MidnightBlue}{-8}{28}{2026}{David Muñoz Gómez}

    \begin{description}
        \item[Asignatura] Álgebra II.
        \item[Curso Académico] 2025/2026.
        \item[Grado] Doble Grado en Matemáticas y Física.
        \item[Grupo] Único.
        \item[Profesor] Aurora Inés del Río Cabeza.
        \item[Descripción] Prueba intermedia I.
        \item[Fecha] 20/04/2026.
        % \item[Duración] ---.
    \end{description}
    \newpage


    % ------------------------------------
    
    \begin{ejercicio}[1 punto]
        En las siguientes cuestiones responde ``VERDADERO'' o ``FALSO'' y haz un breve razonamiento para justificar tu respuesta. Cada respuesta correcta puntúa $0.1$ puntos.

        \begin{enumerate}
            \item Si $G$ es un grupo y para cada elemento $a \in G$ se verifica $a^2 = 1$, entonces $G$ es un grupo abeliano.
            \item El conjunto $\mathbb{R}^* = \mathbb{R} \setminus \{0\}$ con la operación 
            \[
            x \cdot y = \begin{cases} 
            xy & \text{si } x > 0 \\ 
            \nicefrac{x}{y} & \text{si } x < 0 
            \end{cases}
            \]
            es un grupo abeliano.
            \item En un grupo de orden primo cada elemento, distinto del elemento neutro, es un generador.
            \item Un subconjunto de un grupo $G$, que es cerrado para la operación de $G$, no necesariamente es un subgrupo.
            \item Sea $H < G$ un subgrupo y $a \in G$, si $aH = H$, entonces $a \in H$.
            \item Una aplicación inyectiva entre dos grupos es un monomorfismo.
            \item Si $f : G \to H$ es un monomorfismo de grupos y $G$ es abeliano, entonces, $H$ es abeliano.
            \item Sean $H_1, H_2 < G$ subgrupos de índices finitos y $[G : H_1]$ y $[G : H_2]$ son primos relativos, entonces $G = H_1 \vee H_2$.
            \item Dos grupos finitos que tienen subgrupos propios isomorfos son necesariamente isomorfos.
            \item Si en $S_4$ consideramos el subgrupo $H = \langle (1\ 3\ 4) \rangle$, entonces, los conjuntos cocientes de clases laterales por la derecha $Hx$ y clases laterales por la izquierda $xH$ son iguales.
        \end{enumerate}
    \end{ejercicio}

    \begin{ejercicio}[1 punto]
        Considera el grupo $C_8 = \langle a \mid a^8 = 1 \rangle$. Calcula:
        \begin{enumerate}
            \item Todos los automorfismos de $C_8$.
            \item El orden de cada uno de los automorfismos de $C_8$.
            \item Identifica el grupo $\text{Aut}(C_8)$ con un grupo conocido.
        \end{enumerate}
    \end{ejercicio}

    \newpage
    \setcounter{ejercicio}{0} % Reiniciar contador de ejercicios
    \begin{ejercicio}[1 punto]
            En las siguientes cuestiones responde ``VERDADERO'' o ``FALSO'' y haz un breve razonamiento para justificar tu respuesta. Cada respuesta correcta puntúa $0.1$ puntos.
    \begin{enumerate}
            \item Si $G$ es un grupo y para cada elemento $a \in G$ se verifica $a^2 = 1$, entonces $G$ es un grupo abeliano.
            
            \textbf{VERDADERO.} Si $a^2 = 1$ para todo $a \in G$, entonces cada elemento es su propio inverso ($a = a^{-1}$). Dados $a, b \in G$, se tiene:
            \[
            ab = (ab)^{-1} = b^{-1}a^{-1} = ba
            \]
            Por lo tanto, $G$ es abeliano.
            \item El conjunto $\mathbb{R}^* = \mathbb{R} \setminus \{0\}$ con la operación 
            \[
            x \cdot y = \begin{cases} 
            xy & \text{si } x > 0 \\ 
            \nicefrac{x}{y} & \text{si } x < 0 
            \end{cases}
            \]
            es un grupo abeliano.


            \textbf{FALSO.} La operación no es conmutativa. Por ejemplo, si tomamos $x = -1$ e $y = 2$:
            \[
            x \cdot y = \frac{-1}{2} = -\frac{1}{2}
            \]
            Mientras que para $y \cdot x$, como $y = 2 > 0$:
            \[
            y \cdot x = 2 \cdot (-1) = -2
            \]
            Como $x \cdot y \neq y \cdot x$, no puede ser un grupo abeliano.
            \item En un grupo de orden primo cada elemento, distinto del elemento neutro, es un generador.
            
            \textbf{VERDADERO.} Sea $G$ un grupo de orden primo $p$. Por el teorema de Lagrange, el orden de cualquier subgrupo de $G$ debe dividir a $p$. Como $p$ es primo, los únicos subgrupos posibles son el trivial $\{1\}$ y el propio $G$. Si tomamos cualquier elemento $a \in G$ distinto del elemento neutro ($a \neq 1$), el subgrupo generado $\langle a \rangle$ no es el trivial, por lo que obligatoriamente $\langle a \rangle = G$, haciendo de $a$ un generador.
            \item Un subconjunto de un grupo $G$, que es cerrado para la operación de $G$, no necesariamente es un subgrupo.
            
            \textbf{VERDADERO.} Por ejemplo, en el grupo infinito $(\mathbb{Z}, +)$, el subconjunto de los enteros positivos $\mathbb{Z}^+$ es cerrado bajo la suma, pero no contiene el elemento neutro ($0$) ni los elementos inversos, por lo que no es un subgrupo.
            \item Sea $H < G$ un subgrupo y $a \in G$, si $aH = H$, entonces $a \in H$.
            
            \textbf{VERDADERO.} Como $H$ es un subgrupo, contiene al elemento neutro $e \in~H$. Por tanto, $a = ae \in aH$. Dado que por hipótesis $aH = H$, se concluye directamente que $a \in H$.
            \item Una aplicación inyectiva entre dos grupos es un monomorfismo.
            
            \textbf{FALSO.} No es cierto, puesto que puede no respetar la estructura de grupo. Por ejemplo, consideremos los grupos $G,H = (\mathbb{Z}, +)$ y la aplicación
            \Func{f}{G}{H}{n}{n+1}

            Veamos que $f$ es inyectiva:
            \begin{equation*}
            f(n) = f(m) \implies n + 1 = m + 1 \implies n = m
            \end{equation*}

            Sin embargo, $f$ no es un homomorfismo de grupos, puesto que:
            \begin{equation*}
            f(n + m) = (n + m) + 1 \neq (n + 1) + (m + 1) = f(n) + f(m)
            \end{equation*}
             Por lo tanto, $f$ es una aplicación inyectiva que no es un monomorfismo de grupos.
            \item Si $f : G \to H$ es un monomorfismo de grupos y $G$ es abeliano, entonces, $H$ es abeliano.
            
            \textbf{FALSO.} Sea $G = \{1\}$ (que es abeliano) y $H = S_3$ (que es no abeliano). La aplicación trivial $f: G \to H$ definida por $f(1) = \id$ es un homomorfismo inyectivo (monomorfismo), pero $H$ no es abeliano.
            \item Sean $H_1, H_2 < G$ subgrupos de índices finitos y $[G : H_1]$ y $[G : H_2]$ son primos relativos, entonces $G = H_1 \vee H_2$.
            
            \textbf{VERDADERO.} Por las propiedades de los índices de subgrupos, se cumple que $[G : H_1 \cap H_2] \le [G : H_1][G : H_2]$, dándose la igualdad si y solo si $G = H_1 H_2$ (que equivale a que el supremo o unión de subgrupos $H_1 \vee H_2$ sea todo $G$). Como los índices son primos relativos, se alcanza dicha igualdad, por lo que $G = H_1 \vee H_2$.
            \item Dos grupos finitos que tienen subgrupos propios isomorfos son necesariamente isomorfos.
            
            \textbf{FALSO.} Consideremos el grupo cíclico $C_4$ y el grupo de Klein $V_4$. Ambos son grupos de orden 4 no isomorfos entre sí ($C_4 \not\cong V_4$). Sin embargo, ambos poseen subgrupos propios de orden 2 que son isomorfos a $C_2$.
            \item Si en $S_4$ consideramos el subgrupo $H = \langle (1\ 3\ 4) \rangle$, entonces, los conjuntos cocientes de clases laterales por la derecha $Hx$ y clases laterales por la izquierda $xH$ son iguales.
            
            \textbf{FALSO.} Para ver que $xH \neq Hx$, basta comprobar que la conjugación de un elemento de $H$ por $x$ se sale del subgrupo $H$:

            $$x (1\ 3\ 4) x^{-1} = (1\ 2)(1\ 3\ 4)(1\ 2)^{-1} = (1\ 2)(1\ 3\ 4)(1\ 2) = (2\ 3\ 4)$$

            Dado que $(2\ 3\ 4) \notin H$, sabemos automáticamente que $xHx^{-1} \neq H$, y por lo tanto, $xH \neq Hx$. 

            Si calculamos explícitamente ambas clases laterales, la diferencia es evidente:
            \begin{itemize}
                \item \textbf{Clase por la izquierda:} $$xH = (1\ 2)H = \{(1\ 2), (1\ 2)(1\ 3\ 4), (1\ 2)(1\ 4\ 3)\} = \{(1\ 2), (1\ 3\ 4\ 2), (1\ 4\ 3\ 2)\}$$
                \item \textbf{Clase por la derecha:} $$Hx = H(1\ 2) = \{(1\ 2), (1\ 3\ 4)(1\ 2), (1\ 4\ 3)(1\ 2)\} = \{(1\ 2), (1\ 2\ 3\ 4), (1\ 2\ 4\ 3)\}$$
            \end{itemize}
        \end{enumerate}
    \end{ejercicio}

    \begin{ejercicio}[1 punto]
        Considera el grupo $C_8 = \langle a \mid a^8 = 1 \rangle$. Calcula:
        \begin{enumerate}
            \item Todos los automorfismos de $C_8$.
            
            Sabemos que $\operatorname{Aut}(C_8) \cong \mathcal{U}(\mathbb{Z}_8)$ mediante el isomorfismo:
            \[ \varphi: \mathcal{U}(\mathbb{Z}_8) \to \operatorname{Aut}(C_8), \quad k \mapsto f_k \]
            donde el automorfismo $f_k$ actúa elevando a la potencia $k$, es decir, $f_k(x) = x^k$ para todo $x \in C_8$.
            
            Calculamos las unidades de $\mathbb{Z}_8$ (los coprimos con 8 menores que 8): 
            $$\mathcal{U}(\mathbb{Z}_8) = \{1, 3, 5, 7\}$$
            
            Por tanto, hay exactamente 4 automorfismos en $\operatorname{Aut}(C_8)$:
            \begin{itemize}
                \item $f_1(a) = a^1 = a \quad$ (el automorfismo identidad)
                \item $f_3(a) = a^3$
                \item $f_5(a) = a^5$
                \item $f_7(a) = a^7$
            \end{itemize}


            \item El orden de cada uno de los automorfismos de $C_8$.
            
            Gracias al isomorfismo anterior, el orden de cada automorfismo $f_k$ bajo la operación de composición de funciones es idéntico al orden del elemento $k$ en el grupo multiplicativo $\mathcal{U}(\mathbb{Z}_8)$.
        
            Calculamos los órdenes elevando al cuadrado cada elemento en $\mathbb{Z}_8$:
            \begin{itemize}
                \item Para $f_1$: $1^1 \equiv 1 \pmod 8$. Su orden es \textbf{1}.
                \item Para $f_3$: $3^2 = 9 \equiv 1 \pmod 8$. Su orden es \textbf{2}.
                \item Para $f_5$: $5^2 = 25 \equiv 1 \pmod 8$. Su orden es \textbf{2}.
                \item Para $f_7$: $7^2 = 49 \equiv 1 \pmod 8$. Su orden es \textbf{2}.
            \end{itemize}
            \item Identifica el grupo $\text{Aut}(C_8)$ con un grupo conocido.
            
            Como $|\operatorname{Aut}(C_8)| = 4$. Salvo isomorfismo, los únicos grupos de orden 4 son el grupo cíclico $C_4$ y el grupo de Klein $V_4$

            Observando los órdenes calculados en el apartado anterior, vemos que el orden máximo de los elementos de $\operatorname{Aut}(C_8)$ es 2. Al no existir ningún elemento de orden 4, el grupo no puede estar generado por un solo elemento y, por tanto, no es cíclico. 
            
            En conclusión, identificamos el grupo de automorfismos con el producto directo de dos grupos cíclicos de orden 2:
            \[ \operatorname{Aut}(C_8) \cong V_4 \]
        \end{enumerate}
    \end{ejercicio}
\end{document}